Das zusätzliche Gewicht $m\,g$ der Ladung muss durch zusätzlichen Auftrieb getragen werden. Sinkt der Tanker (senkrechte Bordwände, Grundfläche $A$) um $\Delta h$ tiefer ein, verdrängt er zusätzlich das Wasservolumen $A\,\Delta h$:
\begin{align}
 \rho_\text{W}\,A\,\Delta h\,g &= m\,g
\end{align}
Mit $\rho_\text{W} = \rhoWO$ (FoTaBe S.~187) folgt
\begin{align}
 \Delta h &= \dhTF \\
   &= \frac{\mOel}{\rhoW\cdot\Atk} \\
   &= \dhT \approx \result{\dhTP{0}{2}}
\end{align}
Die Dichte des Heizöls spielt keine Rolle; es kommt nur auf die Masse der Ladung an. (Im Meerwasser mit etwas grösserer Dichte sinkt der Tanker ein wenig weniger tief ein.)
